\subsection{The Weil norm of a coupled physical cutoff}
\label{subsec:burnol-coupled-cutoff}
The local vectors of
Proposition~\ref{prop:fd-burnol-weil} admit a quantitative
cutoff calculation for each fixed finite polynomial family.
Both endpoints of the physical polynomial enter the
uncut transform; the endpoint at zero determines the leading
sample after a shorter physical cutoff.

Fix $\gamma\in\Gamma$, put $\rho=1/2-i\gamma$, and retain
Burnol's multiplier
\[
 v(s)=\left(\frac{s}{1-s}\right)^3
       2(2\pi)^{-s}\Gamma(s)\cos\frac{\pi s}{2}.
\]
For a polynomial $P(t)=\sum_{k=0}^K b_kt^k$, define
\[
 H_R=\sum_{k=0}^K b_k\mathcal JX_k^{e^{-R}},
 \qquad
 p_R(x)=R^{-1/2}P(-x/R)e^{i\gamma x}
                    {\bf1}_{[-R,0]}(x).
\]
All constants below may depend on this fixed ordinate and
polynomial. Choose a real nonnegative
$\chi\in C_c^\infty((-1,1))$ equal to one near zero, and set
$\chi_L(x)=\chi(x/L)$.

\begin{proposition}\label{prop:burnol-coupled-cutoff}
Suppose $\int_0^1P(t)\,dt=0$, that $P$ has exact order $M\ge0$
at zero, and that it vanishes to order at least $M$ at one.
Write
\[
 a=\frac{P^{(M)}(0)}{M!}\ne0,\quad
 A_{\chi,M}=\int_0^1\chi(-t)t^Mdt,\quad
 B_{\chi,M}=\int_0^1\chi(-t)^2t^{2M}dt.
\]
Then $q(H_R)=O_{\gamma,P}(R^{-2M-1})$. For every
$L=L(R)\to\infty$ with $L/R\to0$,
\begin{align}
 q(\chi_LH_R)
 &\sim q((1-\chi_L)H_R)
 \sim m_\gamma|a|^2A_{\chi,M}^2
                  \frac{L^{2M+2}}{R^{2M+1}},
                 \label{eq:burnol-cutoff-q}\\
 \|\chi_LH_R\|_2^2
 &\sim |a|^2B_{\chi,M}
                  \frac{L^{2M+1}}{R^{2M+1}}.
                 \label{eq:burnol-cutoff-physical}
\end{align}
In particular,
\[
 \frac{q(\chi_LH_R)}{\|\chi_LH_R\|_2^2}
 \sim m_\gamma\frac{A_{\chi,M}^2}{B_{\chi,M}}\,L.
\]
\end{proposition}

\begin{proof}
We first record the endpoint cancellations in the
multiplier formula. Put $s_\nu=1/2-i\nu$ and
$d_\nu=s_\nu-\rho=-i(\nu-\gamma)$. Summing
\eqref{eq:fd-burnol-transform} against $b_kR^{-k-1/2}$
and applying Leibniz's rule gives, for $\nu\ne\gamma$,
\begin{align}
 \widehat H_R(\nu)
 =R^{-1/2}\bigg\{&
 \sum_{j=0}^K\frac{P^{(j)}(0)}{R^jd_\nu^{j+1}}\notag\\
 &-\frac{e^{-Rd_\nu}}{v(s_\nu)}
 \sum_{j=0}^K\frac{P^{(j)}(1)}{j!R^j}
       \partial_\rho^j\left(\frac{v(\rho)}{d_\nu}\right)
 \bigg\}+\mathcal R_R(\nu).
 \label{eq:burnol-endpoint-transform}
\end{align}
Indeed $\sum_{k\ge j}b_k\binom{k}{j}=P^{(j)}(1)/j!$,
while $b_jj!=P^{(j)}(0)$. The lower-tail remainder bound
\eqref{eq:fd-burnol-cutoff-error} gives
\begin{equation}\label{eq:burnol-poly-rem}
 |\mathcal R_R(\nu)|
 \le C_{\gamma,P}R^{-M-1/2}
           e^{-5R/2}/|s_\nu|.
\end{equation}
Here $b_k=0$ for $k<M$. The two endpoint sums in
\eqref{eq:burnol-endpoint-transform} start at $j=M$.
Isolation of the fixed ordinate and $|v(s_\nu)|=1$ therefore
give
\[
 |\widehat H_R(\nu)|
 \le C_{\gamma,P}R^{-M-1/2}|\nu-\gamma|^{-1}
      +|\mathcal R_R(\nu)|\quad(\nu\in\Gamma,\ \nu\ne\gamma).
\]
The squared majorants are summable against $m_\nu$.

At the target coordinate the direct integral calculation
used for \eqref{eq:fd-burnol-target} instead gives
\[
 \widehat H_R(\gamma)=
 \sqrt R\sum_{j=0}^K
 \frac{v^{(j)}(\rho)}{v(\rho)j!R^j}
             \int_0^1P^{(j)}(t)\,dt
          +O_{\gamma,P}(R^{-M-1/2}).
\]
The $j=0$ term vanishes by the integral constraint. For
$1\le j\le M$, the integral equals
$P^{(j-1)}(1)-P^{(j-1)}(0)=0$. Every remaining term is
$O(R^{-M-1/2})$. This proves $q(H_R)=O(R^{-2M-1})$.
It also identifies the two leading off-target endpoint terms:
for each fixed $\nu\ne\gamma$,
\[
 \widehat H_R(\nu)=
 \frac{R^{-1/2}}{-i(\nu-\gamma)}
 \left\{P(0)-P(1)\frac{v(\rho)}{v(s_\nu)}
                     e^{iR(\nu-\gamma)}\right\}
 +O_{\gamma,P,\nu}(R^{-3/2}).
\]

We next control multiplication by $\chi_L$ in the sampling
norm. Set $w(t)=\log(3+|t|)$ and
$\|F\|_w^2=\int w(t)|F(t)|^2dt$.
The counted unit-interval zero bound implies, for
$0<\epsilon\le1$,
\begin{equation}\label{eq:burnol-short-sampling}
 \sum_{\nu\in\Gamma}m_\nu|F(\nu)|^2
 \le C\left(\epsilon^{-1}\|F\|_w^2
                         +\epsilon\|F'\|_w^2\right).
\end{equation}
To prove it, apply the elementary one-dimensional Sobolev
inequality at $\nu$ on $[\nu-\epsilon,\nu+\epsilon]$.
At a point $t$, at most $C\log(3+|t|)$ counted ordinates
have intervals containing it. Summation and integration
give \eqref{eq:burnol-short-sampling}.

If $e$ has Fourier transform in $L^2(w\,dt)$, multiplication
by $\chi_L$ is convolution with
$k_L(t)=L\widehat\chi(Lt)/(2\pi)$ in frequency.
The inequality
$w(t+u)\le Cw(t)\log(3+|u|)$ and Minkowski's inequality give
\[
 \|k_L*F\|_w\le C_\chi\|F\|_w,\qquad
 \|k_L'*F\|_w\le C_\chi L\|F\|_w
 \qquad(L\ge1).
\]
Both follow by substituting $uL$ in the weighted $L^1$
norms of the Schwartz kernels. Taking $\epsilon=1/L$ in
\eqref{eq:burnol-short-sampling} proves the useful bound
\begin{equation}\label{eq:burnol-cutoff-sampling}
 q(\chi_Le)\le C_\chi L\|\widehat e\|_w^2.
\end{equation}
This explicitly uses weighted Fourier control; a physical
$L^2$ bound alone has not been substituted for a Weil bound.

The exact difference from the polynomial indicator is
\begin{equation}\label{eq:burnol-poly-error}
 \widehat H_R(\nu)-\widehat p_R(\nu)
 =\frac{R^{-1/2}e^{-Rd_\nu}}{v(s_\nu)}
 \sum_{j=M}^K\frac{P^{(j)}(1)}{j!R^j}\,Q_j(\nu)
       +\mathcal R_R(\nu),
\end{equation}
where
\[
 Q_j(\nu)=\partial_\rho^j
             \frac{v(s_\nu)-v(\rho)}{s_\nu-\rho}.
\]
Each singularity at $\gamma$ is removable, and
$Q_j(\nu)=O_{\gamma,j}(|\nu|^{-1})$ at infinity.
Consequently $Q_j\in L^2(w\,d\nu)$. On the real line the
two factors $e^{-Rd_\nu}$ and $v(s_\nu)^{-1}$ have absolute
value one. Thus
\[
 \|\widehat H_R-\widehat p_R\|_w
       \le C_{\gamma,P}R^{-M-1/2}.
\]
Equations \eqref{eq:burnol-cutoff-sampling} and Fourier
Plancherel now give
\begin{equation}\label{eq:burnol-two-cutoff-errors}
 q\bigl(\chi_L(H_R-p_R)\bigr)
       \le C_{\gamma,P,\chi}L R^{-2M-1},\qquad
 \|\chi_L(H_R-p_R)\|_2
       \le C_{\gamma,P}R^{-M-1/2}.
\end{equation}
These statements also place the compactly supported
functions in the completed form domain. For example,
physical mollification supported within the fixed support
margin converges in the weighted Fourier Sobolev norm
appearing in \eqref{eq:burnol-short-sampling}; it gives
compact smooth approximants in the same window.

It remains to compute the compact polynomial itself.
Write $g_{R,L}=\chi_Lp_R$. At the target,
\begin{align*}
 \widehat g_{R,L}(\gamma)
 &=R^{-1/2}L\int_0^1\chi(-t)P(Lt/R)\,dt\\
 &=aA_{\chi,M}\frac{L^{M+1}}{R^{M+1/2}}
                      \left(1+O_P(L/R)\right).
\end{align*}
At a distinct ordinate, integration by parts $M+1$ times
gives
\[
 |\widehat g_{R,L}(\nu)|
 \le C_{P,\chi}R^{-M-1/2}|\nu-\gamma|^{-M-1}
       \qquad(1\le L\le R/2).
\]
For completeness, all left endpoint terms vanish because
$\chi$ has compact support in $(-1,1)$. At zero the
derivatives of order less than $M$ vanish, the derivative
of order $M$ is $(-1)^MP^{(M)}(0)/R^M$, and the
$L^1$ norm of the derivative of order $M+1$ of
$\chi(x/L)P(-x/R)$ on $[-L,0]$ is $O_{P,\chi}(R^{-M})$.
The last bound follows term by term from Leibniz's rule,
using $P^{(j)}(t)=O_P(t^{M-j})$ for $j\le M$ and $L/R\le1/2$.
Therefore the off-target squared samples sum to
$O_{\gamma,P,\chi}(R^{-2M-1})$, whereas the target contributes
the leading expression in \eqref{eq:burnol-cutoff-q}.
The first error norm in \eqref{eq:burnol-two-cutoff-errors}
is smaller than this leading norm by
$O(L^{-M-1/2})$. This proves the asymptotic for
$q(\chi_LH_R)$.
The same leading sample has the exact global jet image
given by \eqref{eq:hb-evaluation-jet}:
\[
 (\mathcal U(\chi_LH_R))^{(m_\gamma-1)}(\gamma)
 =\frac{iF^{(m_\gamma)}(\gamma)}{\sqrt\pi}
                      \widehat{\chi_LH_R}(\gamma).
\]
Thus the coordinate responsible for the cutoff norm is
identified with its corresponding derivative evaluation,
with the multiplicity and normalization retained.

Direct substitution also gives
\[
 \|g_{R,L}\|_2^2
 =R^{-1}L\int_0^1\chi(-t)^2|P(Lt/R)|^2dt
 \sim |a|^2B_{\chi,M}\frac{L^{2M+1}}{R^{2M+1}}.
\]
The second error bound in \eqref{eq:burnol-two-cutoff-errors}
is smaller than its norm by $O(L^{-M-1/2})$, proving
\eqref{eq:burnol-cutoff-physical}. Finally
$\|H_R\|_q=O(R^{-M-1/2})$ is smaller than
$\|\chi_LH_R\|_q$ by $O(L^{-M-1})$.
The reverse triangle inequality proves that
$q((1-\chi_L)H_R)\sim q(\chi_LH_R)$.
\end{proof}

\paragraph{An exponential estimate for the local multiplier error.}
The multiplier gives a stronger version of the first
bound in \eqref{eq:burnol-two-cutoff-errors}. With the same
fixed data, for $1\le L\le R/2$,
\begin{equation}\label{eq:burnol-exp-local-error}
 \|\chi_L(H_R-p_R)\|_q
 \le C_{\gamma,P,\chi}R^{-M-1/2}
 \left\{L^{3/2}e^{-2(R-L)}
                    +\sqrt L\,e^{-5R/2}\right\}.
\end{equation}
Here and below $\|f\|_q=q(f)^{1/2}$.
We give the multiplier justification to keep the direction
of the physical translation explicit.

Let $\kappa_j$ be the inverse Fourier transform of
$Q_j(\nu)/v(s_\nu)$. For $x>0$,
\begin{equation}\label{eq:burnol-positive-kernel-decay}
 |\kappa_j(x)|+|\kappa_j'(x)|
                   \le C_{\gamma,j}e^{-2x}.
\end{equation}
Indeed,
\[
 \frac{Q_j(\nu)}{v(s_\nu)}
 =\frac{j!}{d_\nu^{j+1}}
  -\frac1{v(s_\nu)}
        \partial_\rho^j\left(\frac{v(\rho)}{d_\nu}\right).
\]
The apparent real pole cancels in the whole expression.
The reciprocal multiplier is holomorphic in
$0\le\Im\nu<5/2$: its apparent issue at $s=1$ is a zero
of order two, and its first pole above the critical line
is at $s=3$, corresponding to $\nu=5i/2$.
The gamma formula and Stirling's formula give
$|v(5/2-it)^{-1}|\le C(1+|t|)^{-2}$.
Shift the inverse Fourier contour to $\nu=t+2i$.
The vertical edge terms vanish because the whole integrand
before the Fourier exponential is $O(1/|t|)$ uniformly
across this fixed strip. On the shifted line the inverse
transform of $j!/d_\nu^{j+1}$ is zero for $x>0$.
One can check this without a contour closure from
\[
 \frac{j!}{(2-i(t-\gamma))^{j+1}}
     =\int_0^\infty u^je^{-2u}e^{i(t-\gamma)u}\,du:
\]
its inverse transform is supported on $x\le0$.
The remaining shifted-line integral has integrand
$O((1+|t|)^{-3})$, and its derivative in $x$ has an
integrable $O((1+|t|)^{-2})$ majorant. Its factor $e^{-2x}$
proves \eqref{eq:burnol-positive-kernel-decay}, with a
continuous first derivative for $x>0$.

The nonremainder part of \eqref{eq:burnol-poly-error}
therefore has inverse transform
\[
 R^{-1/2}e^{-iR\gamma}
       \sum_{j=M}^K\frac{P^{(j)}(1)}{j!R^j}
                                     \kappa_j(x+R).
\]
On $|x|\le L$ its value and first derivative are bounded
by $CR^{-M-1/2}e^{-2(R-L)}$.
After multiplication by $\chi_L$ its physical $H^1$ norm
is at most $C\sqrt L R^{-M-1/2}e^{-2(R-L)}$.
The compact-support sampling estimate
$\|f\|_q\le CL\|f\|_{H^1}$ follows from
\eqref{eq:burnol-short-sampling} with $\epsilon=1$,
Fourier Plancherel, and
$w(\nu)\le C(1+\nu^2)$ applied to $f$ and $xf$.
This proves the first term in
\eqref{eq:burnol-exp-local-error}. The remainder
\eqref{eq:burnol-poly-rem} has weighted Fourier norm at
most $CR^{-M-1/2}e^{-5R/2}$; applying
\eqref{eq:burnol-cutoff-sampling} proves its second term.

\paragraph{The two-jet family and the all-window lower rate.}
For $P(t)=1-2t$, $M=0$ and $a=1$. Thus, at
$L=c\log\log R$ with any fixed $c>0$,
\[
 q((1-\chi_L)H_R)\sim
 m_\gamma A_{\chi,0}^2
             \frac{c^2(\log\log R)^2}{R},
 \qquad
 \|\chi_LH_R\|_2^2\sim
 B_{\chi,0}\frac{c\log\log R}{R}.
\]
The same asymptotics hold for the exact target-canceling
combination
\[
 Z_R=\mathcal JX_0^{e^{-R}}-c_R\mathcal JX_1^{e^{-R}},
 \qquad
 c_R=
 \frac{\widehat{\mathcal JX_0^{e^{-R}}}(\gamma)}
      {\widehat{\mathcal JX_1^{e^{-R}}}(\gamma)}
     =2+O_\gamma(R^{-1}).
\]
In fact $q(Z_R)=O_\gamma(R^{-1})$, and all $M=0$ cutoff
error estimates above are uniform over bounded polynomial
coefficients. Its local polynomial
$1-c_Rt$ has the same nonzero constant term and the same
leading cutoff samples and physical mass. No exact
integral constraint is needed for that cutoff calculation;
the uncut bound follows explicitly from
\[
 \|Z_R\|_q\le\|H_R\|_q+|c_R-2|\,
                   \|\mathcal JX_1^{e^{-R}}\|_q
                =O_\gamma(R^{-1/2}),
\]
using Proposition~\ref{prop:fd-burnol-weil}.

The proved all-window estimate of
Theorem~\ref{thm:cosine-lower-rate} is
\[
 a_1(L)\ge
 \exp[-C e^{2L/d_{\rm MT}}L(1+\log L)].
\]
At $L=c\log\log R$ its right side is
\[
 \exp\!\left[-C(\log R)^{2c/d_{\rm MT}}\,
        c\log\log R\bigl(1+\log(c\log\log R)\bigr)\right].
\]
The Rayleigh quotient of the compact two-jet
function is asymptotic to
\[
 m_\gamma\frac{A_{\chi,0}^2}{B_{\chi,0}}\,c\log\log R.
\]
The cutoff Weil norm is asymptotic to a positive constant times
$(\log\log R)/\sqrt R$, whereas the uncut Weil norm is
$O_\gamma(R^{-1/2})$.
These formulas quantify the missing term for this fixed
family and window scale. They provide no bound uniform in
the ordinate or as the number of jet orders grows.
The original all-height integer-deficit estimate remains
unproved.
